\documentclass[11pt,openany]{book}

\usepackage[a5paper,top=17mm,bottom=20mm,inner=19mm,outer=15mm,headheight=14pt]{geometry}
\usepackage{fontspec}
\usepackage[vietnamese]{babel}
\usepackage{unicode-math}
\usepackage{microtype}
\usepackage{xcolor}
\usepackage{hyperref}
\usepackage{bookmark}
\usepackage{fancyhdr}
\usepackage{titlesec}
\usepackage[most]{tcolorbox}
\usepackage{fvextra}
\usepackage{listings}
\lstdefinelanguage{text}{}
\usepackage[
  fencedCode,
  pipeTables,
  smartEllipses,
  strikeThrough,
  texMathDollars,
  texMathSingleBackslash,
  hybrid
]{markdown}

\setmainfont{Charis SIL}
\setsansfont{Inter}
\setmonofont{Cousine}[Scale=MatchLowercase]
\setmathfont{Latin Modern Math}
\newfontfamily\starfont{DejaVu Sans}
\providecommand{\square}{\mdlgwhtsquare}

\definecolor{bookblue}{HTML}{164E63}
\definecolor{bookteal}{HTML}{0F766E}
\definecolor{bookgold}{HTML}{B7791F}
\definecolor{bookpaper}{HTML}{F8FAFC}
\definecolor{bookink}{HTML}{17202A}

\hypersetup{
  colorlinks=true,
  linkcolor=bookteal,
  urlcolor=bookblue,
  pdftitle={Đi tìm thử thách 2 — Bản dịch tiếng Việt},
  pdfauthor={Bản dịch tiếng Việt dùng cho mục đích giáo dục}
}
\DeclareRobustCommand{\difficulty}[1]{{\starfont\color{bookgold}%
  \ifcase#1\or ★\or ★★\or ★★★\or ★★★★\fi}}
\pdfstringdefDisableCommands{%
  \def\difficulty#1{[#1 sao]}%
}

\color{bookink}
\setlength{\parindent}{1.1em}
\setlength{\parskip}{0.15em}
\setlength{\emergencystretch}{2em}
\setlength{\headheight}{15pt}
\raggedbottom
\renewcommand{\chaptername}{Bài}

\titleformat{\chapter}[display]
  {\sffamily\bfseries\color{bookblue}}
  {\filleft\Large BÀI \thechapter}
  {0.4ex}
  {\titlerule[1.2pt]\vspace{1ex}\Huge}
\titleformat{\section}
  {\sffamily\Large\bfseries\color{bookteal}}
  {\thesection}{0.65em}{}
\titleformat{\subsection}
  {\sffamily\large\bfseries\color{bookblue}}
  {\thesubsection}{0.55em}{}
\titleformat{\subsubsection}
  {\sffamily\normalsize\bfseries\color{bookblue}}
  {\thesubsubsection}{0.5em}{}

\pagestyle{fancy}
\fancyhf{}
\fancyhead[LE]{\sffamily\small\color{bookblue}\leftmark}
\fancyhead[RO]{\sffamily\small\color{bookblue}\rightmark}
\fancyfoot[C]{\sffamily\small\thepage}
\renewcommand{\headrulewidth}{0.3pt}

\renewenvironment{quote}
  {\begin{tcolorbox}[enhanced,breakable,colback=bookpaper,colframe=bookteal,
    boxrule=0pt,leftrule=2pt,arc=0pt,left=8pt,right=8pt,top=5pt,bottom=5pt]}
  {\end{tcolorbox}}

\DefineVerbatimEnvironment{Highlighting}{Verbatim}{
  breaklines,breakanywhere,fontsize=\small,
  frame=single,framesep=5pt,rulecolor=\color{bookteal}}

% Pseudocode is visually distinct from sample input/output blocks.
\newtcolorbox{pseudocodebox}{
  enhanced,breakable,
  colback=bookblue!3,colframe=bookteal!75,
  colbacktitle=bookteal!12,coltitle=bookblue,
  title={Pseudocode},fonttitle=\sffamily\bfseries\footnotesize,
  boxrule=.55pt,arc=1.5mm,
  left=5pt,right=5pt,top=4pt,bottom=4pt,
  before skip=7pt,after skip=7pt}

\def\markdownRendererInputFencedCode#1#2{%
  \ifstrequal{#2}{pseudocode}{%
    \begin{pseudocodebox}
      \VerbatimInput[fontsize=\scriptsize,breaklines,breakanywhere,
        breaksymbolleft={},breaksymbolright={}]{#1}%
    \end{pseudocodebox}%
  }{%
    \VerbatimInput[fontsize=\small,breaklines,breakanywhere]{#1}%
  }%
}

% Map Markdown headings to the book hierarchy: # = chapter, ## = section, ...
\def\markdownRendererHeadingOne#1{\chapter{#1}}
\def\markdownRendererHeadingTwo#1{\section{#1}}
\def\markdownRendererHeadingThree#1{\subsection{#1}}
\def\markdownRendererHeadingFour#1{\subsubsection{#1}}
\def\markdownRendererLineBreak{\newline}
\def\markdownRendererImage#1#2#3#4{%
  \begin{figure}[htbp]
    \centering
    \includegraphics[width=\linewidth,height=.43\textheight,keepaspectratio]{#3}%
    \caption{#1}%
  \end{figure}}

\begin{document}

\frontmatter
\begin{titlepage}
  \pagecolor{bookpaper}
  \centering
  \vspace*{18mm}
  {\sffamily\bfseries\color{bookblue}\fontsize{29}{34}\selectfont
    ĐI TÌM THỬ THÁCH 2\par}
  \vspace{5mm}
  {\color{bookteal}\rule{0.72\textwidth}{1.4pt}\par}
  \vspace{7mm}
  {\sffamily\Large Các bài toán từ Kỳ thi Lập trình\par
    Sinh viên Ba Lan 2011--2014\par}
  \vfill
  \begin{tcolorbox}[width=.82\textwidth,colback=white,colframe=bookteal,
    boxrule=.7pt,arc=2mm]
    \centering\sffamily
    \textbf{\ifdefined\FullBook Bản dịch tiếng Việt\else Bản xem trước tiếng Việt\fi}\par
    \smallskip
    \ifdefined\FullBook 44 bài toán, 2011--2014\else 4 bài đầu tiên của năm 2011\fi
  \end{tcolorbox}
  \vfill
  {\small\itshape Bản dịch phục vụ mục đích giáo dục.\par}
  \vspace{10mm}
  {\sffamily\color{bookgold}BẢN THẢO • 2026}
\end{titlepage}
\nopagecolor

\chapter*{Ghi chú về bản dịch}
\addcontentsline{toc}{chapter}{Ghi chú về bản dịch}
\ifdefined\FullBook
Bản PDF này là bản thảo tiếng Việt gồm đủ 44 bài toán và lời giải.
\else
Bản PDF này minh họa cách trình bày dự kiến cho bản dịch hoàn chỉnh.
\fi
Nội dung
được ưu tiên dịch từ bản tiếng Anh. Với 33 lời giải chưa có trong bản tiếng Anh,
bản tiếng Ba Lan năm 2019 (phiên bản 2.1) sẽ là nguồn chính.

Tên biến, công thức, độ phức tạp và dữ liệu vào/ra được giữ nguyên về ngữ nghĩa.
Các chương trong bản xem trước vẫn cần một vòng soát kỹ thuật độc lập trước khi
được xem là hoàn chỉnh.

\tableofcontents

\mainmatter
\ifdefined\FullBook
\markdownInput{build/markdown/2011-ary-hinh-chu-nhat-cap-so-cong.md}
\markdownInput{build/markdown/2011-baj-cuoc-dua-duong-pho-bytean.md}
\markdownInput{build/markdown/2011-czy-lieu-no-co-dung.md}
\markdownInput{build/markdown/2011-drz-cay-va-dan-kien.md}
\markdownInput{build/markdown/2011-eks-diet-soc-dat.md}
\markdownInput{build/markdown/2011-fra-may-giat.md}
\markdownInput{build/markdown/2011-gen-bo-sinh-bit.md}
\markdownInput{build/markdown/2011-her-tra-chieu.md}
\markdownInput{build/markdown/2011-ilo-chi-so-thong-minh.md}
\markdownInput{build/markdown/2011-jas-hang-dong.md}
\markdownInput{build/markdown/2011-krz-nhen-chu-thap.md}
\markdownInput{build/markdown/2012-aut-may-ban-hang-tu-dong.md}
\markdownInput{build/markdown/2012-biu-chuyen-xe-buyt.md}
\markdownInput{build/markdown/2012-cia-day-so.md}
\markdownInput{build/markdown/2012-dna.md}
\markdownInput{build/markdown/2012-ewa-tinh-gia-tri-bieu-thuc.md}
\markdownInput{build/markdown/2012-for-cong-thuc-mot.md}
\markdownInput{build/markdown/2012-gen-cuu-khung-long.md}
\markdownInput{build/markdown/2012-hyd-hydra.md}
\markdownInput{build/markdown/2012-inw-nghich-the.md}
\markdownInput{build/markdown/2012-jut-de-mai-lam.md}
\markdownInput{build/markdown/2012-kro-tho.md}
\markdownInput{build/markdown/2013-aut-duong-cao-toc.md}
\markdownInput{build/markdown/2013-baj-bytehattan.md}
\markdownInput{build/markdown/2013-cie-tho-moc.md}
\markdownInput{build/markdown/2013-dem-bieu-tinh.md}
\markdownInput{build/markdown/2013-egz-ky-thi.md}
\markdownInput{build/markdown/2013-fot-may-ban-toc-do.md}
\markdownInput{build/markdown/2013-gra-tro-choi-bi.md}
\markdownInput{build/markdown/2013-her-nguoi-hung.md}
\markdownInput{build/markdown/2013-inz-ky-thuat-di-truyen.md}
\markdownInput{build/markdown/2013-jan-janosik.md}
\markdownInput{build/markdown/2013-koc-chan.md}
\markdownInput{build/markdown/2014-adw-luat-su.md}
\markdownInput{build/markdown/2014-ben-xang.md}
\markdownInput{build/markdown/2014-cen-gia-ca.md}
\markdownInput{build/markdown/2014-dzi-uoc-so.md}
\markdownInput{build/markdown/2014-euk-nim-euclid.md}
\markdownInput{build/markdown/2014-fil-cot-tru.md}
\markdownInput{build/markdown/2014-glo-nong-len-toan-cau.md}
\markdownInput{build/markdown/2014-hit-ban-hit-cua-mua.md}
\markdownInput{build/markdown/2014-ins-dan-dung.md}
\markdownInput{build/markdown/2014-jas-hang-dong.md}
\markdownInput{build/markdown/2014-kap-thuyen-truong.md}

\else
\markdownInput{build/markdown/2011-ary-hinh-chu-nhat-cap-so-cong.md}
\markdownInput{build/markdown/2011-baj-cuoc-dua-duong-pho-bytean.md}
\markdownInput{build/markdown/2011-czy-lieu-no-co-dung.md}
\markdownInput{build/markdown/2011-drz-cay-va-dan-kien.md}
\fi

\backmatter
\chapter*{Thông tin nguồn}
\addcontentsline{toc}{chapter}{Thông tin nguồn}
\begin{itemize}
  \item \emph{Looking for a Challenge 2}, bản tiếng Anh, Warszawa, 2019.
  \item \emph{W poszukiwaniu wyzwań 2}, bản tiếng Ba Lan, ấn bản sửa đổi lần hai,
        phiên bản 2.1, Warszawa, 2019.
\end{itemize}

\end{document}
